\section{System design}

\subsection{Architecture overview}
Figure~\ref{fig:architecture} shows the whole system, from the two sources to the three ways of serving
results. Every box runs in its own Docker container, and one Docker Compose file starts all of them.

\begin{figure}[!htb]
\centering
% FleetPulse architecture (Lambda). Drawn at text width (16 cm), no scaling needed.
\begin{tikzpicture}[
  font=\sffamily\scriptsize,
  box/.style={draw=black!55, rounded corners=2pt, align=center, inner sep=2.5pt, line width=0.5pt},
  arrow/.style={-{Stealth[length=2mm]}, line width=0.7pt, draw=black!70},
  dashedarrow/.style={arrow, dashed},
  lbl/.style={font=\sffamily\tiny, text=black, inner sep=2.4pt},
  layer/.style={font=\sffamily\scriptsize\bfseries, text=black},
]
% ---- layer titles
\node[layer] at (1.25, 8.00) {SOURCES};
\node[layer] at (4.60, 8.00) {INGESTION};
\node[layer] at (8.30, 8.00) {PROCESSING};
\node[layer] at (11.90, 8.00) {STORAGE};
\node[layer] at (14.90, 8.00) {SERVING};

% ---- sources (x 0.10 to 2.30)
\node[box, fill=black!6, text width=2.2cm, minimum height=1.4cm] (ss) at (1.25, 6.40)
  {\textbf{Stream simulator}\\Python, 30 vehicles,\\one ping per second};
\node[box, fill=black!6, text width=2.2cm, minimum height=1.4cm] (bs) at (1.25, 2.50)
  {\textbf{Batch simulator}\\Python, cost file\\once per day};

% ---- ingestion (x 3.20 to 5.90)
\node[box, fill=orange!14, text width=2.75cm, minimum height=2.3cm] (kafka) at (4.60, 6.40)
  {\textbf{Apache Kafka 4.3}\\KRaft, one broker\\[2pt]topic\\\texttt{fleet.telemetry.raw}\\6 partitions,\\key = vehicle id\\[2pt]+ dead-letter topic};
\node[box, fill=orange!14, text width=2.75cm, minimum height=1.4cm] (land) at (4.60, 2.50)
  {\textbf{Landing zone}\\daily CSV file and\\manifest (SHA-256)};

% ---- processing (x 6.75 to 9.85)
\node[box, fill=teal!14, text width=2.9cm, minimum height=2.3cm] (speed) at (8.30, 6.40)
  {\textbf{Speed layer}\\Spark 4.1 Structured\\Streaming: validate,\\de-duplicate, watermark,\\zone join, 1 h and 1 day\\windows, idle alerts};
\node[box, fill=violet!12, text width=2.9cm, minimum height=0.9cm] (lake) at (8.30, 4.45)
  {\textbf{Master dataset}\\immutable raw Parquet};
\node[box, fill=blue!10, text width=2.9cm, minimum height=1.9cm] (batch) at (8.30, 2.50)
  {\textbf{Batch layer}\\Airflow 3.3 DAG runs a\\Spark 4.1 batch job\\once per day: exact\\recompute and join\\with partner costs};

% ---- storage (x 10.75 to 13.05)
\node[box, fill=violet!12, text width=2.1cm, minimum height=5.4cm] (pg) at (11.90, 4.45)
  {\textbf{PostgreSQL 17}\\[5pt]\texttt{speed.*}\\real-time views\\[5pt]\texttt{batch.*}\\exact daily views\\[5pt]\texttt{serving.*}\\merge views\\[5pt]\texttt{ops.*}\\audit, alerts, clock};

% ---- serving (x 13.80 to 16.00)
\node[box, fill=green!12, text width=2.0cm, minimum height=1.0cm] (api) at (14.90, 6.40)
  {\textbf{FastAPI}\\REST and /metrics};
\node[box, fill=green!12, text width=2.0cm, minimum height=1.0cm] (graf) at (14.90, 4.45)
  {\textbf{Grafana 12}\\3 dashboards};
\node[box, fill=green!12, text width=2.0cm, minimum height=1.0cm] (rep) at (14.90, 2.50)
  {\textbf{Daily report}\\HTML and CSV};

% ---- data flow
\draw[arrow] (ss) -- node[lbl, above] {JSON} (kafka);
\draw[arrow] (kafka) -- node[lbl, above] {stream} (speed);
\draw[arrow] (bs) -- node[lbl, above] {CSV} (land);
\draw[arrow] (land) -- node[lbl, above] {sensor} (batch);
\draw[dashedarrow] (speed) -- node[lbl, right] {raw archive} (lake);
\draw[arrow] (lake) -- node[lbl, right] {one day} (batch);
\draw[arrow] (speed.east) -- node[lbl, above] {upsert} (speed.east -| pg.west);
\draw[arrow] (batch.east) -- node[lbl, above] {one txn} (batch.east -| pg.west);
\draw[arrow] (pg.east |- api.west) -- (api.west);
\draw[arrow] (pg.east |- graf.west) -- (graf.west);
\draw[arrow] (pg.east |- rep.west) -- (rep.west);

% ---- observability band (x 0.10 to 16.00)
\node[box, fill=yellow!16, text width=15.6cm, minimum height=0.8cm, anchor=west] (obs) at (0.10, 0.60)
  {\textbf{Observability:} JSON logs from every stage \,\textbullet\, Prometheus 3 metrics (producers, Spark listener,
   API, Airflow via StatsD) \,\textbullet\, 14 alert rules $\rightarrow$ Alertmanager $\rightarrow$ API webhook
   \,\textbullet\, lineage and audit tables};
\foreach \n in {bs, land, batch, pg, rep} {
  \draw[densely dotted, draw=black!45, line width=0.6pt] (\n.south) -- (\n.south |- obs.north);
}
\end{tikzpicture}

\caption{FleetPulse architecture. The top row is the streaming path and the bottom row is the daily batch
path. Both meet in PostgreSQL. The dotted lines show that every stage sends logs and metrics.}
\label{fig:architecture}
\end{figure}

\textbf{Streaming path.} The stream simulator sends a JSON event to Kafka for every vehicle every real
second. Spark Structured Streaming reads the topic with four queries. Three of them clean the events and
update the real-time tables in PostgreSQL. The fourth writes every raw record, unchanged, to the master
dataset in Parquet.

\textbf{Batch path.} Once per simulated day the batch simulator drops a CSV file and a manifest into the
landing zone. An Airflow sensor finds the file, checks it and copies it into the master dataset. Airflow
then starts a Spark batch job. The job reads one full day of raw events, removes bad and repeated records,
joins the day with the cost file and writes exact results to PostgreSQL in a single transaction.

\textbf{Serving.} PostgreSQL views merge the two layers. The FastAPI service, the Grafana dashboards and the
daily report all read from these tables and views.

\subsection{Data sources}
The streaming source sends one JSON event per Kafka message with the fields \code{event\_id} (a UUID, also
sent as the Kafka header \code{trace\_id}), \code{event\_type} (\code{ping} every 2 simulated minutes,
\code{trip\_start} or \code{trip\_end}), \code{event\_ts} (simulated time), \code{vehicle\_id},
\code{driver\_id}, \code{trip\_id} (empty when idle), \code{status} (\code{idle}, \code{enroute} to a
pickup or \code{on\_trip}), \code{lat}, \code{lon}, \code{speed} (km/h), \code{odometer\_km} and
\code{fare} (the running fare, or the final fare on \code{trip\_end}). A \code{trip\_end} event also carries
the pickup point, trip distance and duration. We added \code{event\_id}, \code{event\_type} and
\code{odometer\_km} to the fields listed in the brief.

The daily source is a CSV file from the fuel and garage partners, with a manifest. It has one row per
vehicle and day with the fields \code{record\_id}, \code{expense\_date}, \code{vehicle\_id},
\code{partner\_id}, \code{fuel\_litres}, \code{fuel\_cost}, \code{maintenance\_cost} (upkeep, wear and any
service bill), \code{distance\_covered} (kilometres reported by the partner, which we check against GPS) and
\code{service\_flag} (\code{Y} on a day with a garage service).

To test the cleaning logic, both sources add problems on purpose: about 1\% of events are sent twice,
0.5\% of pings are broken (bad JSON, missing ids, a GPS point at 0,0, negative speed and more), and
vehicles sometimes lose signal and send their events late. Most daily files also contain a repeated row, a
row for an unknown vehicle, a bad number or a missing vehicle id.

\subsection{Kafka design}
We use two topics. \code{fleet.telemetry.raw} has six partitions and carries all telemetry, which the speed
layer reads and archives for the batch layer. \code{fleet.telemetry.dlq} has one partition. It is the
dead-letter queue, which keeps invalid events with the reason for checks and replay. Messages are keyed by
\code{vehicle\_id}, so all events of one vehicle go to the same partition and stay in order. This order
matters for idle streaks and trips. Six partitions allow up to six consumers to read in parallel. The
producer uses \code{acks=all} with the idempotent producer switched on, so retries after a network error never
create duplicates in Kafka. It also uses lz4 compression and a 20~ms linger time to send small batches.

\subsection{Storage design}
The \textbf{master dataset} is a folder of Parquet files that we never change, only add to. Raw telemetry is
partitioned by the simulated date (\code{lake/raw/telemetry/event\_date=2026-09-01/}) and written by the
Spark file sink, whose metadata log makes each write exactly-once. The partner files are kept in
\code{lake/raw/expenses/}. The batch layer also keeps its own results in \code{lake/curated/} and reads the
last two days from there to find profit trends.

\textbf{PostgreSQL 17} is the serving store. We chose it because the views are small (a few hundred rows per
day), and it gives SQL, joins, upserts and transactions. It has five schemas: \code{ref} (vehicles and zones),
\code{speed} (live zone and vehicle metrics, vehicle state and idle alerts), \code{batch} (exact metrics,
trips, profitability, cost rates, quarantined cost rows and the daily summary), \code{serving} (views that
merge the two layers and the live profit projection) and \code{ops} (simulated clock, micro-batch ledger,
dead letters, run audit trail, alert notifications and reconciliation).
